\documentclass[10pt,a4paper]{amsart}
\usepackage[margin=19mm]{geometry}
\usepackage{amsmath,amssymb,mathtools}
\usepackage[scr=rsfs,cal=euler]{mathalfa}
\usepackage{xfrac}
\usepackage[unicode,hidelinks,bookmarks=false]{hyperref}
\newcommand{\ie}{i.e.\ }
\newcommand{\cf}{cf.\ }
\newcommand{\resp}{respectively\ }
\newcommand{\Subsection}{\subsection}
\providecommand{\nobreakdash}{\nobreak\hbox{-}}
\setlength{\emergencystretch}{2em}
\multlinegap0pt
\usepackage{mathtools}
\usepackage{amssymb}
\usepackage{dsfont}
\usepackage{tcolorbox}
\newtheorem{lemma}{Lemma}
\usepackage{cleveref}
\crefname{lemma}{Lemma}{Lemmas}
\newcommand{\J}{\mathcal{J}}
\newcommand{\R}{\mathbb{R}}
\newcommand{\U}{\mathcal{U}}
\newcommand{\de}{\,\mathrm{d}}
\title{Ensemble Optimal Control for managing drug resistance in cancer therapies}
\author{Alessandro Scagliotti, Federico Scagliotti, Laura Deborah Locati, Federico Sottotetti}
\date{Source: arXiv:2503.08927v3 (CC BY 4.0)}
\begin{document}
\maketitle

\noindent\emph{Source and licence.} Ensemble Optimal Control for managing drug resistance in cancer therapies. Version arXiv:2503.08927v3 (CC BY 4.0), distributed by arXiv under the Creative Commons Attribution 4.0 International licence, http://creativecommons.org/licenses/by/4.0/, as recorded in the arXiv licence metadata for this exact version and shown on the abstract page https://arxiv.org/abs/2503.08927v3. Full licence notice: \texttt{licenses/arxiv-2503.08927-CC-BY-4.0.txt}.\par\smallskip

\noindent\textit{Editorial scope.} Counted unit: the article's lemma lem:extension (source0.tex lines 488--491). The article's complete original proof of it is delivered with it (source0.tex lines 492--495, one \texttt{proof} environment of 4 lines). No mathematics is written, repaired or re-derived: the statement and the proof are the article's own lines, in the article's own notation, and no retained line is substituted or re-flowed.  Retained uncounted as the source-local context the proof needs: lines 339--341; lines 344--346; lines 349--352; lines 360--364; lines 373--376; lines 471--474.  Nothing was omitted from any retained span, so the package carries no omission record.  No retained line contains a citation, so the wrapper carries no bibliography and no bibliography entry.  No retained line contains a figure environment, an image-inclusion command or drawing code, so no image is delivered and the package declares no asset dependency.  Editorial formatting only, in addition: an AMS \texttt{amsart} preamble that restates the article's own declarations and macros used by the retained lines, namely the environments \texttt{lemma} on separate counters, together with the standard packages those lines need.  The retained environments print in this document as Lemma 1 in order of appearance, whereas the article prints the same items with the numbering of its own full document.  The complete native source of the article as distributed in the arXiv e-print for this exact version is bundled unchanged as \texttt{sources/174607/source0.tex}.
for every $\theta \in \Theta$, and we introduce \begin{equation} \label{eq:simplex_start}
     \Delta \coloneqq \left\{ X \in \R^n : X_i \geq 0 \ \forall i, \, 0\leq \sum_{i=1}^n X_i \leq K_{\max} \right\},
\end{equation} 

\begin{equation}\label{eq:def_adm_ctrls}
    \U\coloneqq \{ {D} \in L^2([0,T],\R): 0\leq {D(t)\leq D_{\max} } \, \mbox{ for a.e. } t \}.
\end{equation}

\begin{lemma}  \label{lem:domain_invariance}
    Let us consider ${D}\in \U$ defined as in \cref{eq:def_adm_ctrls}, and the simplex $\Delta \in \R^{n}$ introduced in \cref{eq:simplex_start}.
    For every $\theta \in \Theta$, if ${X_D^\theta(0)}\in \Delta^{\theta}$, then ${X_D^\theta(t)} \in \Delta^{\theta}$ {(and, in particular, $X_D^\theta(t) \in \Delta$)} for every $t \in [0,T]$.
\end{lemma}

\begin{equation} \label{eq:def_trunc_fields}
    F_0^\theta({X}) \coloneqq \rho({X})\tilde F_0^\theta ({X}) 
    \quad \mbox{and} \quad
    F_1^\theta ({X}) \coloneqq \rho({X}) \tilde F_1^\theta ({X}) 
\end{equation}

\begin{equation} \label{eq:def_ens_funct}
    \J({D}) \coloneqq  \int_\Theta  \int_0^T \ell\left({N_D}^\theta({t} )\right)  
    \de {t} \de\mu(\theta),
\end{equation}

\begin{equation} \label{eq:def_ens_funct'}
    \J'({D})\coloneqq \int_\Theta  \int_0^T \ell\left({N_D^\theta(t)}\right)  
    \de {t} \de\mu(\theta),
\end{equation}

\begin{lemma} \label{lem:extension}
    Let $\J\colon \U\to\R$ and $\J'\colon L^2([0,T],\R)\to \R$ be defined as in \cref{eq:def_ens_funct} and \cref{eq:def_ens_funct'}, respectively.
    If ${X_D^\theta}(0) = \left( {f_1N_0,\ldots, f_nN_0} \right)\in \Delta^{\theta}$ for every $\theta \in \Theta$, then $\J({D})=\J'({D})$ for every ${D} \in\U$.
\end{lemma}

\begin{proof}
    If ${D}\in\U$ and ${X_D^\theta}(0) = \left( {f_1N_0,\ldots, f_nN_0} \right)\in \Delta^{\theta}$ for every $\theta \in \Theta$, then \cref{lem:domain_invariance} implies that ${X_D^\theta(t)} \in \Delta^{\theta} {\subset \Delta}$ for every ${t}\in[0,T]$.
    Recalling that, for every $\theta\in\Theta$, $F_0^\theta \equiv \tilde F_0^\theta$ and $F_1^\theta \equiv \tilde F_1^\theta$ on $B_{2{K_{\max}}}(0)\supset \Delta$ (see \cref{eq:def_trunc_fields}), it turns out that $\J({D})=\J'({D})$.
\end{proof}

\end{document}
